\section{Preliminaries: Hoeffding's lemma}
\label{sec:hoeffding-lemma}

Throughout, $c, C, C_1, C_2, \ldots$ denote universal positive constants whose
values may change from line to line and depend only on quantities specified in
each statement's hypotheses. We write $\Pr[\cdot]$ for probability and
$\E[\cdot]$ for expectation. For a real random variable $Y$, its
\emph{moment generating function} (MGF) is $\lambda \mapsto \E[e^{\lambda Y}]$,
defined for all $\lambda \in \R$ whenever $Y$ is bounded.

The single auxiliary ingredient is Hoeffding's lemma: a bounded, centered
random variable has a sub-Gaussian MGF with the sharp variance proxy
$(b-a)^2/4$. The leading constant $2$ in the final tail bound is a direct
consequence of the constant $1/8$ below.

\begin{lemma}[Hoeffding's lemma]\label{lem:hoeffding}
Let $Y$ be a random variable with $\E[Y] = 0$ and $Y \in [a, b]$ almost surely,
where $a < b$. Then for every $\lambda \in \R$,
\begin{align}\label{eq:hoeffding-lemma}
\E\!\left[ e^{\lambda Y} \right] \;\le\; \exp\!\left( \frac{\lambda^2 (b - a)^2}{8} \right).
\end{align}
\end{lemma}

\begin{proof}
Fix $\lambda \in \R$ and let $\psi(\lambda) := \log \E[e^{\lambda Y}]$ be the
cumulant generating function of $Y$; since $Y$ is bounded, $\E[e^{\lambda Y}]$ is
finite and $\psi$ is infinitely differentiable on $\R$, with derivatives
obtained by differentiating under the expectation. We compute the first two
derivatives and bound the second.

\smallskip\noindent\textbf{Step 1 (boundary values).}
\begin{align}
\psi(0) &\;=\; \log \E[e^{0}] \;=\; \log 1 \;=\; 0, \label{eq:psi-zero} \\
\psi'(\lambda) &\;=\; \frac{\E[Y e^{\lambda Y}]}{\E[e^{\lambda Y}]},
  \qquad\text{so}\qquad \psi'(0) \;=\; \frac{\E[Y]}{\E[1]} \;=\; 0, \label{eq:psi-prime}
\end{align}
where the first line uses $\E[e^0]=1$, and the second differentiates
$\psi=\log\E[e^{\lambda Y}]$ and then sets $\lambda=0$ using $\E[Y]=0$.

\smallskip\noindent\textbf{Step 2 (the second derivative is a variance).}
Define, for each fixed $\lambda$, the tilted law $Q_\lambda$ under which $Y$ has
density proportional to $e^{\lambda Y}$ with respect to its original law. Then
\begin{align}
\psi''(\lambda)
&\;=\; \frac{\E[Y^2 e^{\lambda Y}]}{\E[e^{\lambda Y}]}
       - \left( \frac{\E[Y e^{\lambda Y}]}{\E[e^{\lambda Y}]} \right)^{\!2} \notag \\
&\;=\; \E_{Q_\lambda}[Y^2] - \big( \E_{Q_\lambda}[Y] \big)^2
   \;=\; \Var_{Q_\lambda}(Y), \label{eq:psi-second}
\end{align}
where the first step is the quotient rule applied to $\psi'$ in
Eq.~\eqref{eq:psi-prime}, and the second step rewrites each ratio as an
expectation under $Q_\lambda$ and recognises the difference as a variance.

\smallskip\noindent\textbf{Step 3 (variance of a bounded variable).}
Under $Q_\lambda$ the variable $Y$ still takes values in $[a,b]$, since tilting
only reweights the original law and does not enlarge the support. Writing
$m := (a+b)/2$ for the midpoint, the variance is controlled by the second moment
about $m$:
\begin{align}
\Var_{Q_\lambda}(Y)
&\;=\; \E_{Q_\lambda}\!\big[(Y - \E_{Q_\lambda} Y)^2\big]
   \;\le\; \E_{Q_\lambda}\!\big[(Y - m)^2\big] \notag \\
&\;\le\; \left( \frac{b - a}{2} \right)^{\!2}
   \;=\; \frac{(b - a)^2}{4}, \label{eq:var-bound}
\end{align}
where the first step is the definition of variance, the second uses that the
mean minimises the mean-square deviation (so replacing $\E_{Q_\lambda} Y$ by the
fixed $m$ only increases the value), and the third uses the pointwise bound
$|Y - m| \le (b-a)/2$, valid because $Y \in [a,b]$.

\smallskip\noindent\textbf{Step 4 (Taylor expansion).}
By Taylor's theorem with Lagrange remainder, there exists $\xi$ between $0$ and
$\lambda$ such that
\begin{align}
\psi(\lambda)
&\;=\; \psi(0) + \lambda\,\psi'(0) + \frac{\lambda^2}{2}\,\psi''(\xi) \notag \\
&\;=\; \frac{\lambda^2}{2}\,\psi''(\xi)
   \;\le\; \frac{\lambda^2}{2} \cdot \frac{(b - a)^2}{4}
   \;=\; \frac{\lambda^2 (b - a)^2}{8}, \label{eq:psi-taylor}
\end{align}
where the first step is the second-order Taylor expansion of $\psi$ about $0$,
the second substitutes $\psi(0) = \psi'(0) = 0$ from
Eq.~\eqref{eq:psi-zero} and Eq.~\eqref{eq:psi-prime}, and the inequality combines
Eq.~\eqref{eq:psi-second} with the variance bound in
Eq.~\eqref{eq:var-bound}. Exponentiating both sides yields
Eq.~\eqref{eq:hoeffding-lemma}. This completes the proof.
\end{proof}
